%% ============================================================
%% AULA 10 — Aprender sem professor: autoencoders e embeddings
%% GBC073 — Inteligência Computacional (FACOM/UFU)
%%
%% Esqueleto com esboço de conteúdo — a desenvolver.
%% Ementa (ficha SEI 5116656): item 2 — aprendizagem
%% não-supervisionada; aplicações (clusterização).
%%
%% Compilar:  latexmk -xelatex aula10.tex   (duas passadas!)
%% ============================================================
\documentclass[aspectratio=169, 11pt]{beamer}

\makeatletter
\def\input@path{{./}{../}}
\makeatother

\usetheme{InteligenciaComputacional}   % ANTES do babel: fontspec vem primeiro
\usepackage{babel}
\babelprovide[import, main]{brazilian}

\title[Aula 10 — Autoencoders]{Aula 10: Aprender sem professor — autoencoders e embeddings}
\subtitle[GBC073 — Inteligência Computacional]{Autoencoders \textbullet\ embeddings \textbullet\ aprendizado de representações}
\author[Prof.\ Marcelo Albertini]{Prof.\ Marcelo Keese Albertini}
\institute{GBC073 — Inteligência Computacional \textbullet\ Universidade Federal de Uberlândia}
\date{2026}

\begin{document}

% ------------------------------------------------------------ capa
\begin{frame}[plain, noframenumbering]
  \titlepage
\end{frame}

% ------------------------------------------------------------ roteiro
\begin{frame}{Roteiro da aula}
  \begin{itemize}\setlength{\itemsep}{0.5ex}
    \item<1-> \textbf{Aprender sem rótulos: o que é a tarefa}
    \item<2-> \textbf{Autoencoders (reconstrução, denoising, esparso)}
    \item<3-> \textbf{Embeddings e aprendizado de representações}
    \item<4-> \textbf{Não-supervisionado como pré-treino}
  \end{itemize}
\end{frame}

% ------------------------------------------------------------
\section{Aprender sem rótulos}

\begin{frame}{Aprender sem rótulos: o que é a tarefa}
  \begin{itemize}\setlength{\itemsep}{0.4ex}
    \item Rótulos são caros; dados não rotulados são abundantes.
    \item Sem alvo, o sinal de aprendizado precisa vir dos próprios dados.
    \item A tarefa muda: reconstruir, prever, comparar — a \emph{representação}
          vira o produto.
    \item A ficha da disciplina: ``aprendizagem não-supervisionada'' — aqui,
          com o olhar de deep learning.
  \end{itemize}
\end{frame}

% ------------------------------------------------------------
\section{Autoencoders}

\begin{frame}{Autoencoders (reconstrução, denoising, esparso)}
  \begin{itemize}\setlength{\itemsep}{0.4ex}
    \item Codificador comprime $\vx$ em um código $\vet{z}$; decodificador
          reconstrói $\hat{\vx}$.
    \item A perda é a reconstrução — nenhum rótulo envolvido.
    \item O gargalo força o código a capturar o essencial dos dados.
    \item Variantes: denoising (reconstruir do ruído) e esparso (poucos
          neurônios ativos).
  \end{itemize}
  \begin{ideiabox}
    Um autoencoder é um professor que se ensina a si mesmo: a prova é
    reconstruir o que viu — e a cola é entender a estrutura dos dados.
  \end{ideiabox}
\end{frame}

% ------------------------------------------------------------
\section{Embeddings}

\begin{frame}{Embeddings e aprendizado de representações}
  \begin{itemize}\setlength{\itemsep}{0.4ex}
    \item \emph{Embedding}: o vetor que representa um objeto em um espaço
          aprendido.
    \item Distâncias no espaço de embedding refletem semelhança semântica.
    \item Redução de dimensionalidade como leitura: PCA, t-SNE, UMAP.
    \item A representação aprendida transfere para tarefas supervisionadas.
  \end{itemize}
\end{frame}

% ------------------------------------------------------------
\section{Não-supervisionado como pré-treino}

\begin{frame}{Não-supervisionado como pré-treino}
  \begin{itemize}\setlength{\itemsep}{0.4ex}
    \item Pré-treino não-supervisionado + ajuste fino com poucos rótulos.
    \item O modelo de linguagem prevê o próximo token — a maior tarefa
          não-supervisionada já construída.
    \item Onde mais: representações para busca, recomendação e dados
          científicos.
    \item Não-supervisionado não substitui o supervisionado — o alimenta.
  \end{itemize}
\end{frame}

\end{document}
